\part{Part VII: 结论与展望}
\chapterimage{chapter_head_2.pdf}

\chapter{研究结论}

\section{主要发现}
本研究围绕小红书笔记热度的影响因素与发布前预测展开分析，形成了从数据清洗、探索性分析、特征构造到分类和回归建模的完整流程。三个业务数据压缩包共包含323,598条原始记录，经统一去重和有效性检查后保留199,473条笔记。研究使用“点赞数+收藏数+$4\times$评论数”衡量互动热度，并以10项不包含发布后互动的结构化特征完成两类预测任务。

探索性分析与模型结果共同表明，发布者粉丝基础是最稳定的热度信号；内容长度、情感表达、标签、发布时间和位置元数据提供了较弱但可利用的增量信息。部分变量与热度之间存在非线性关系，因此树模型整体上比线性模型具有更好的排序能力。不过，这些关系来自观察数据，只能解释为统计关联和预测信号，不能直接推断为因果效应。

分类任务中，梯度提升树取得最高PR-AUC（0.440）和ROC-AUC（0.736），适合对潜在高热度笔记排序；线性SVM取得最高召回率（0.651），适合低成本的广覆盖初筛；逻辑回归取得最高默认阈值F1（0.414），可作为易解释基准。回归任务中，梯度提升树取得最低RMSLE（1.849）、最高Spearman相关系数（0.515）和0.429的Top 20\%召回率。接近0的原始尺度$R^2$同时说明，当前模型能够改善相对排序，却难以准确预测极端爆款的绝对互动量。

\section{实践含义}
在现有证据范围内，模型更适合充当内容决策的辅助排序工具，而不是自动化承诺工具。对于候选名额有限的运营场景，可使用梯度提升树按预测分数输出Top-$K$；对于不希望漏掉潜在爆款的初筛场景，可使用类别加权的线性模型，并由人工或第二阶段模型处理较多的假阳性。任何阈值都应在独立验证集上根据误报与漏报成本确定，不能机械采用0.5。

特征重要性可帮助识别值得进一步分析的方向，但不应被转写为“增加正文长度必然提高热度”等确定性建议。更稳妥的做法是将模型发现用于提出运营假设，再通过随机实验、准实验或严格的时间外回测验证其效果。

\chapter{研究局限与未来工作}

\section{研究局限}
第一，本文虽已完成一次时间外切分和一次博主完全隔离的分组切分，但尚未使用多个滚动时间窗口检验平台规则与季节变化。第二，输入主要是文本表层和账号结构化特征，没有稳定配对的图片、完整文本语义、曝光量、传播网络与早期互动序列。第三，样本集中于美妆和家电领域，跨领域泛化能力尚未验证。第四，互动热度是研究构造指标，评论权重取4具有业务解释，但仍需通过敏感性分析检验不同权重下结论是否稳定。

此外，极端爆款受到平台临时流量、突发热点和站外传播影响，这些信息在当前数据中不可观测。因而，模型对数量级和排序的改善不能消除社交媒体传播本身的高度不确定性。

\section{未来工作}
近期工作应把现有单次时间外验证扩展为多个滚动窗口，并继续完成阈值优化、概率校准和跨领域性能比较。对于回归任务，可进一步比较分位数回归、Tweedie损失或两阶段模型，为长尾热度提供区间预测。

在数据条件允许后，可使用中文预训练语言模型提取标题与正文语义，使用视觉编码器提取封面和配图特征，并在稳定笔记标识下进行晚期融合。若能够获得严格按时间排序的发布者历史和用户---笔记交互边，可进一步研究时序模型与图网络。所有新增模型都应在相同的数据切分和无泄漏约束下与当前结构化基准比较，并保留可执行Notebook、环境说明和机器可读结果。

\section{总结}
本研究结果表明，有限的发布前结构化特征已经能够为小红书笔记热度提供优于随机基线的排序信息，但对极端传播结果的解释能力仍然有限。相比追求未经验证的复杂模型，更可靠的路径是先建立可复现基准、明确使用边界，再通过时间外验证和新增模态逐步扩展。由此形成的研究框架既能够支持当前的内容筛选任务，也为后续多模态和动态传播研究保留了清晰接口。
